\documentclass[12pt]{article}
% Préambule commun à tous les documents extraits de « La Revue CAPES »
\usepackage[T1]{fontenc}
\usepackage[utf8]{inputenc}
\usepackage{lmodern}
\usepackage{amsmath,amssymb,amsthm,amsfonts,mathrsfs}
\usepackage{setspace}
\usepackage{listings}
\usepackage{pgfplots}
\pgfplotsset{compat=1.18}
\usepackage{longtable}
\usepackage{adjustbox}
\usepackage{lettrine}
\usepackage{tkz-tab}
\usepackage{changepage}
\usepackage{pdflscape}
\usepackage{graphicx}
\usepackage{tikz}
\usepackage{xcolor}
\usepackage[a4paper,top=2cm,bottom=2.5cm,left=2.5cm,right=2cm]{geometry}
\usepackage{fancyhdr}
\usepackage{draftwatermark}
\SetWatermarkText{La RevueCapes}
\SetWatermarkLightness{0.9}
\SetWatermarkScale{2}
\usepackage{hyperref}
\hypersetup{colorlinks=true,linkcolor=blue!50!black,urlcolor=blue!50!black}

\definecolor{revue}{RGB}{20,60,120}

\pagestyle{fancy}
\fancyhf{}
\renewcommand{\headrulewidth}{0.4pt}
\setlength{\headheight}{15pt}

% \entete{Matière}{Titre}{Type de document}
\newcommand{\entete}[3]{%
  \fancyhead[L]{\small\textsc{La Revue CAPES} -- #1}%
  \fancyhead[R]{\small #2}%
  \fancyfoot[C]{\thepage}%
  \thispagestyle{plain}%
  \begin{center}
    {\color{revue}\rule{\textwidth}{1.2pt}}\\[4pt]
    {\small\textsc{Concours directs CAP-CEG / CAPES -- Burkina Faso}}\\[6pt]
    {\LARGE\bfseries\color{revue} #2}\\[6pt]
    {\large\itshape #1 -- #3}\\[2pt]
    {\color{revue}\rule{\textwidth}{1.2pt}}
  \end{center}
  \vspace{0.5cm}%
}

% Encadré pour les remarques ajoutées dans les corrigés
\newcommand{\remarque}[1]{\par\medskip\noindent\fcolorbox{revue}{revue!6}{\parbox{\dimexpr\linewidth-2\fboxsep-2\fboxrule}{\textbf{\color{revue}Remarque.} #1}}\par\medskip}


\begin{document}
\entete{Mathématiques}{Sujet type n°2}{Énoncé}
\begin{spacing}{1.5}

\textbf{Exercice 1}

 Soient les réels $a,b$ et $c$.On considère les matrices carrées $A$ et $B$ suivantes:

$
A=\begin{pmatrix}
a+b & b+c & c+a \\ 
a^2+b^2 & b^2+c^2 & c^2+a^2 \\ 
a^3+b^3 & b^3+c^3 & c^3+a^3
\end{pmatrix}$ et $B=\begin{pmatrix}
1 & 1 & 1 \\ 
c & a & b \\ 
c^2 & a^2 & b^2
\end{pmatrix} $

(a) Calculer le déterminant de $B$ en fonction de $a,b$ et $c$.

(b) Montrer que $\det(A)=2abc\times\det(B)$.

(c) Calculer $det(A)$ en fonction de $a,b$ et $c$.

(d) En déduire le rang de $A$ suivant les valeurs de $a,b$ et $c$.

\textbf{Exercice 2 }

 Soit $E$ un $\mathbb{R}$-espace vectoriel de dimension 3 et $B=\langle u,v,w\rangle$ une base de E.
 On considère l'endomorphisme $f$ de $E$ défini par
  $f(u)=u+2v-w, f(v)=3u-w, f(w)=-3u+2v+3w$.
 
 1. Donner la matrice $A$ de $f$  relativement à la base $B$.
 
 2. Déterminer le spectre de $A$.
 
 3. Déterminer les sous-espaces propres de $A$.
 
 4. Montrer que la matrice $A$ est diagonalisable?
 
 5. Calculer $A^n$ pour tout $n\in \mathbb{N}$.
 
 6. Quel est l'image du vecteur $u+v+w$,par l'endomorphisme $f\circ f\circ f$?
 
 7. On considère trois suites réelles $(a_n)_{n\in \mathbb{N}},(b_n)_{n\in \mathbb{N}},(c_n)_{n\in \mathbb{N}}$ définies par les relations de récurrence
 
\[
\begin{cases}
a_{n+1} = a_n+3b_n-3c_n\\
b_{n+1} = 2a_n+2c_n\\
c_{n+1} = -a_n-b_n+3c_n
\end{cases}
\]
et $a_0=b_0=c_0=1$.

Déterminer $a_n,b_n,c_n$ en fonction de $n$.
 

\medskip\textbf{Exercice 3}
 
  1) Dresser un tableau donnant tous les résultats possibles de lancer de 2 dés équilibrés à 6 faces.
 
 La variable aléatoire $X$ désigne le résultat du premier dé.
 
 La variable aléatoire $Y$ désigne le résultat du deuxième dé.
 
 On considère que les variables aléatoires $X$ et $Y$ sont indépendantes.
 
 2) Établir la loi de probabilité de la variable aléatoire somme $S=X+Y$, donnant la somme des résultats des 2 dés.
 
 \medskip\textbf{Exercice 4}
 
 1. Soit le système 
 
 $X'=AX$ avec $A=\begin{pmatrix}
 1 & -2 & 0 \\ 
 2 & 0 & -1 \\ 
 4 & -2 & -1
 \end{pmatrix} $
 
 Déterminer les valeurs propres de $A$ puis un système de solutions de $X'=AX$.
 
 2. Même exercice avec
 $A=\begin{pmatrix}
 1 & -1 \\ 
 4 & -3
 \end{pmatrix} $, puis $A=\begin{pmatrix}
 2 & 1 & 0 & 0 \\ 
 0 & 0 & -1 & 0 \\ 
 1 & 2 & 1 & 0 \\ 
 -1 & 2 & 2 & 1
 \end{pmatrix} $
 
 
 \quad
 

\end{spacing}
\end{document}
